\section{Introduction}

A normal office network never carries just one protocol. At any moment there is
web traffic, IoT sensor data, video calls and background sync jobs all sharing
the same wire. Studying how these protocols behave together, how they fail, and
how easy they are to tell apart from the outside, needs a testbed that can
actually produce that mix on demand and let you turn the dials while it runs.

This project builds exactly that testbed, based on the assignment in
\cite{hsrwassignment}. It generates traffic for five protocols at the same
time: HTTP/2 \cite{rfc7540}, QUIC/HTTP-3 \cite{rfc9000, rfc9114}, MQTT
\cite{mqtt5} with all three QoS levels, and raw TCP and UDP. Everything runs
as 15 Docker containers orchestrated with Docker Compose \cite{dockercompose},
started with one \texttt{docker-compose up}. A central controller loads a YAML profile
and tells four traffic generators what to send and how fast; a web dashboard
lets you change any of that live, without restarting anything.

Beyond the basic rate and payload-size knobs, the system implements four
sending patterns (constant, periodic burst, linear ramp and Poisson-distributed
random intervals), a stealth mode that makes the raw TCP/UDP generator
statistically indistinguishable from normal traffic, and an autonomous Adaptive
Control loop that scales a generator's rate up or down based on the error rate
it is currently seeing, with no human in the loop.

The rest of this report is organized as follows. Section II describes the
architecture. Section III walks through the implementation decisions that
mattered, and why they were made that way. Section IV documents the
configuration format. Sections V through VII present the Wireshark-based
traffic analysis, the failure-visibility results and the single- versus
multi-machine comparison. Section VIII reflects on what the data actually
showed about protocol behavior, and Section IX documents how AI tools were
used while building this.
